\chapter{Quick start}
\label{ch:quickstart}

This chapter takes you from nothing to audio going through a plugin on the
server, in about ten minutes.

\section{Get the software}

\fxme\ is built from source for now. You need a C++17 compiler, CMake 3.22 or
newer, and the JUCE framework. From the \fxme\ folder:

\begin{lstlisting}
cmake -S . -B build
cmake --build build --target FxmeNetServer -j2
cmake --build build --target FxmeNetPlugin_VST3 -j2
cmake --build build --target FxmeNetTop -j2
\end{lstlisting}

This gives three programs: the server (\cmd{fxmenet-server}), the plugin
(\file{FxmeNet.vst3}) and the monitoring tool (\cmd{fxmenet-top},
chapter~\ref{ch:top}). \todo{where the build leaves the files; how JUCE and
FxmeTools are found (sibling clones or pinned versions)}

\section{Install the plugins on the server}

The server offers the VST3 effects it finds in its plugin folders: the
standard VST3 folders of the machine (on Linux, typically
\file{~/.vst3}, \file{/usr/lib/vst3} and
\file{/usr/local/lib/vst3}), plus any you add (section~\ref{sec:folders}).
Install the plugins you want to use there, the way you would on any computer.

\section{Start the server}

On the server machine, in a terminal:

\begin{lstlisting}
fxmenet-server
\end{lstlisting}

The first time, the server examines every plugin it finds (this can take a
minute for a large collection) and then logs the list, ending with lines like:

\begin{lstlisting}[language={}]
2026-11-10 14:02:11  48 plugins (48 examined, 0 from the cache, 0 failed)
2026-11-10 14:02:11  listening: control TCP 54322, audio UDP 54321, discovery UDP 54320
\end{lstlisting}

Later starts are quick, because only new or changed plugins are examined
again. Leave the server running. Stop it with \key{Ctrl}+\key{C}.

\begin{note}
Your computers' firewall must let the three ports through on the local
network (chapter~\ref{ch:network}). Under Ubuntu's \cmd{ufw}, for example,
allow 54320 to 54322.
\end{note}

\section{Install the plugin on the client}

Copy \file{FxmeNet.vst3} to your VST3 folder on the computer that runs the
DAW (on Windows, \file{C:\Program Files\Common Files\VST3}), then let your DAW rescan its plugins.
On macOS the plugin is also built as an Audio Unit. The AU is not notarized, so macOS may refuse to load it at first; the \file{README} on the project's GitHub page explains what to do. The AU has not been tested yet, and feedback is welcome.

\section{Your first chain}

\begin{enumerate}
  \item In the DAW, add \ui{FxmeNet} to a track, as an effect.
  \item Click \ui{Find}. The server answers on the local network; choose it
    in the list. (If it does not appear, type its address in the field, for
    example \cmd{192.168.1.20}.)
  \item Click \ui{Connect}. The status line next to the button says when
    the session is ready.
  \item The list \ui{SERVER PLUGINS} shows what the server offers. Type a few
    letters of a plugin's name in the \ui{Filter} field, and press
    \key{Return} (or double-click a plugin) to add it to the chain.
  \item Play. The audio goes through the server. The plugin's controls
    appear in the right-hand panel; its parameters are also listed in the
    DAW as automation parameters.
\end{enumerate}

The line at the bottom of the editor shows the latency that \fxme\ reports to
the DAW and the count of concealed blocks. On a healthy network the count
stays at zero, or nearly.

\begin{tip}
To see what the server is doing, open \cmd{fxmenet-top} in a terminal on the
server (chapter~\ref{ch:top}). It shows each session, the server's load and
the traffic.
\end{tip}

\section{Where to go next}

\begin{itemize}
  \item Choosing the block size and budget, and making the server
    reliable: chapter~\ref{ch:performance}.
  \item Parallel chains, sidechains and multichannel plugins:
    chapter~\ref{ch:routing}.
  \item Something does not work: chapter~\ref{ch:troubleshooting}.
\end{itemize}
